\documentclass[letterpaper]{article} % DO NOT CHANGE THIS
\usepackage[submission]{aaai2027}  % DO NOT CHANGE THIS
% The serif, sans-serif, and monospaced fonts are loaded automatically by
% aaai2027.sty (newtxtext, helvet, courier). DO NOT add \usepackage{times},
% \usepackage{helvet}, \usepackage{courier}, or any other font package.
\usepackage[hyphens]{url}  % DO NOT CHANGE THIS
\usepackage{graphicx} % DO NOT CHANGE THIS
\urlstyle{rm} % DO NOT CHANGE THIS
\def\UrlFont{\rm}  % DO NOT CHANGE THIS
\usepackage{natbib}  % DO NOT CHANGE THIS AND DO NOT ADD ANY OPTIONS TO IT
\usepackage{caption} % DO NOT CHANGE THIS AND DO NOT ADD ANY OPTIONS TO IT
\frenchspacing  % DO NOT CHANGE THIS
%
% These are recommended to typeset algorithms but not required. See the subsubsection on algorithms. Remove them if you don't have algorithms in your paper.
% \usepackage{algorithm}
% \usepackage{algorithmic}

%
% These are recommended to typeset listings but not required. See the subsubsection on listing. Remove this block if you don't have listings in your paper.
% \usepackage{newfloat}
% \usepackage{listings}
% \DeclareCaptionStyle{ruled}{labelfont=normalfont,labelsep=colon,strut=off} % DO NOT CHANGE THIS
% \lstset{%
% 	basicstyle={\footnotesize\ttfamily},% footnotesize acceptable for monospace
% 	numbers=left,numberstyle=\footnotesize,xleftmargin=2em,% show line numbers, remove this entire line if you don't want the numbers.
% 	aboveskip=0pt,belowskip=0pt,%
% 	showstringspaces=false,tabsize=2,breaklines=true}
% \floatstyle{ruled}
% \newfloat{listing}{tb}{lst}{}
% \floatname{listing}{Listing}

%
% Recommended for better-looking tables
\usepackage{booktabs}
\usepackage{amsmath,amssymb}

%
% Keep the \pdfinfo as shown here. There's no need
% for you to add the /Title and /Author tags.
\pdfinfo{
/TemplateVersion (2027.1)
}

% DISALLOWED PACKAGES
% \usepackage{authblk} -- This package is specifically forbidden
% \usepackage{balance} -- This package is specifically forbidden
% \usepackage{CJK} -- This package is specifically forbidden
% \usepackage{float} -- This package is specifically forbidden
% \usepackage{flushend} -- This package is specifically forbidden
% \usepackage{fullpage} -- This package is specifically forbidden
% \usepackage{geometry} -- This package is specifically forbidden
% \usepackage{hyperref} -- This package is specifically forbidden
% \usepackage{navigator} -- This package is specifically forbidden
% (or any other package that embeds links such as navigator or hyperref)
% \usepackage{indentfirst} -- This package is specifically forbidden
% \usepackage{layout} -- This package is specifically forbidden
% \usepackage{multicol} -- This package is specifically forbidden
% \usepackage{nameref} -- This package is specifically forbidden
% \usepackage{savetrees} -- This package is specifically forbidden
% \usepackage{setspace} -- This package is specifically forbidden
% \usepackage{stfloats} -- This package is specifically forbidden
% \usepackage{tabu} -- This package is specifically forbidden
% \usepackage{titlesec} -- This package is specifically forbidden
% \usepackage{tocbibind} -- This package is specifically forbidden
% \usepackage{ulem} -- This package is specifically forbidden
% \usepackage{wrapfig} -- This package is specifically forbidden
% DISALLOWED COMMANDS
% \nocopyright -- Your paper will not be published if you use this command
% \addtolength -- This command may not be used
% \balance -- This command may not be used
% \baselinestretch -- Your paper will not be published if you use this command
% \clearpage -- No page breaks of any kind may be used for the final version of your paper
% \columnsep -- This command may not be used
% \newpage -- No page breaks of any kind may be used for the final version of your paper
% \pagebreak -- No page breaks of any kind may be used for the final version of your paper
% \pagestyle -- This command may not be used
% \tiny -- This is not an acceptable font size.
% \vspace{- -- No negative value may be used in proximity of a caption, figure, table, section, subsection, subsubsection, or reference
% \vskip{- -- No negative value may be used to alter spacing above or below a caption, figure, table, section, subsection, subsubsection, or reference

\setcounter{secnumdepth}{0} %May be changed to 1 or 2 if section numbers are desired.

% The file aaai2027.sty is the style file for AAAI Press
% proceedings, working notes, and technical reports.
%

% Title

% Your title must be in mixed case, not sentence case.
% That means all verbs (including short verbs like be, is, using,and go),
% nouns, adverbs, adjectives should be capitalized, including both words in hyphenated terms, while
% articles, conjunctions, and prepositions are lower case unless they
% directly follow a colon or long dash
\title{CaRePrompt: From Cellular Evidence to Hierarchical Region Prompting for Interactive Whole-Slide Tissue Segmentation}
\author{Anonymous Submission}
\affiliations{}

% Example, Single Author, ->> remove \iffalse,\fi and place them surrounding AAAI title to use it
\iffalse
\title{My Publication Title --- Single Author}
\author{
    Author Name
}
\affiliations{
    Affiliation\\
    Affiliation Line 2\\
    name@example.com
}
\fi

% Example, Multiple Authors, ->> remove \iffalse,\fi and place them surrounding AAAI title to use it
\iffalse
\title{My Publication Title --- Multiple Authors}
\author{
    % Authors
    First Author Name\textsuperscript{\rm 1,\rm 2}\equalcontrib,
    Second Author Name\textsuperscript{\rm 2}\equalcontrib,
    Third Author Name\textsuperscript{\rm 1}\corresponding
}
\affiliations{
    % Affiliations
    \textsuperscript{\rm 1}Affiliation 1\\
    \textsuperscript{\rm 2}Affiliation 2\\
    firstAuthor@affiliation1.com, secondAuthor@affilation2.com, thirdAuthor@affiliation1.com
}
\fi

\begin{document}

\maketitle

\begin{abstract}
Interactive tissue segmentation in whole-slide images (WSIs) requires models to identify user-specified tissues from sparse visual prompts while integrating information across cellular, regional, and global tissue scales. Existing approaches often rely on fixed patches or handcrafted regions, which may limit their ability to adapt to complex tissue boundaries and capture biological context across different scales.

We present CaRePrompt, a cell-aware hierarchical prompting framework for interactive target-versus-rest tissue segmentation in WSIs. CaRePrompt learns adaptive tissue region representations from multiscale slide features, incorporates cellular information to enhance region-level understanding, and leverages hierarchical contextual reasoning to connect local morphology with broader tissue organization. By conditioning segmentation on positive and negative visual prompts, CaRePrompt enables a unified model to segment user-specified tissues without requiring class-specific fine-tuning.

We evaluate CaRePrompt on patient-separated WSI cohorts under episodic interactive segmentation settings. Experimental results demonstrate that the proposed framework effectively integrates cellular and spatial context, achieving accurate tissue delineation across diverse tissue targets. These findings highlight the potential of cell-aware hierarchical prompting as a general framework for interactive computational pathology.
\end{abstract}

% Uncomment the following to link to your code, datasets, an extended version or similar.
% You must keep this block between (not within) the abstract and the main body of the paper.
% Make sure that you do not de-anonymize yourself with these links.
% \begin{links}
%     \link{Code}{https://aaai.org/example/code}
%     \link{Datasets}{https://aaai.org/example/datasets}
%     \link{Extended version}{https://aaai.org/example/extended-version}
% \end{links}

\section{Introduction}

CaRePrompt moves interactive WSI tissue segmentation from prompt matching on isolated pixels or patches to prompt-conditioned reasoning over a cell-aware hierarchy of learned tissue regions. This change addresses a limitation that remains after a user has specified the target: the model must still determine which local structures support that target, how those structures compose into tissue, and where the resulting boundary lies at slide scale. On 4,000 held-out, patient-separated test episodes spanning 12 tissue targets and point, small-region, and large-region interactions, a single CaRePrompt model obtains 0.7392 macro Dice and 0.8128 micro Dice without per-class fine-tuning. These results establish the central contribution first: visual prompts can define a reusable target-versus-rest tissue task when cellular, regional, and slide context are represented within the same prediction pathway.

The corresponding failure occurs at the prompt-to-mask interface. A fine patch can preserve nuclei and narrow boundaries but omit the tissue organization needed to interpret them; a coarse patch can capture that organization while averaging away the cellular evidence that distinguishes visually similar tissues. Fixed patches or precomputed superpixels make this trade-off before seeing the prompt. Even when retrieval identifies relevant regions, converting similarity ranks into a mask introduces a second failure point: classwise thresholds, area priors, or top-percentile rules can change the predicted tissue more than the prompt itself. CaRePrompt shifts both decisions into the model. Differentiable region assignments determine the prediction units, cell attention supplies microscopic evidence, sparse parent--child edges provide multiscale context, and a learned decoder maps signed prompt evidence directly to calibrated region probabilities.

Prior work has solved important parts of this problem, but not their coupling. SAM established promptable segmentation and zero-shot transfer in natural images \cite{kirillov2023segment}, while MedSAM adapted box-prompted segmentation to diverse medical modalities \cite{ma2024medsam}. WSI-SAM further introduced multi-resolution features for prompted histopathology segmentation \cite{liu2024wsisam}. Separately, HIPT showed that the natural WSI pyramid supports hierarchical representation learning \cite{chen2022hipt}, and CellViT demonstrated that transformer-based cellular representations can transfer across histopathology datasets \cite{horst2024cellvit}. These advances provide prompts, multiresolution context, or cellular representations, respectively. The missing step is to make them interact at the same learned tissue regions: cellular evidence should alter region semantics, hierarchy should connect those regions across physical scales, and positive and negative prompts should jointly define the binary tissue query rather than select a fixed class head.

Our contributions follow this missing computation and are paired with explicit experimental checks:

\noindent\textbf{Adaptive region-to-pixel prediction.}
We couple differentiable tissue regions with a context-aware region-to-pixel decoder, replacing retrieval cutoffs with soft probability projection. In a controlled 4,000-episode validation comparison under dynamic online region targets, this change raises pixel macro Dice from 0.7212 to 0.7316 and Boundary F1 from 0.2287 to 0.2878.

\noindent\textbf{Cell-aware signed prompting.}
We introduce cell-aware region tokens and a signed set encoder that represents multiple positive and negative prompts without reducing each set to a fixed mean prototype. Controlled 360-episode ablations isolate the two effects: adding the cell branch changes macro/micro Dice by $+0.00013/-0.00007$, while adapting the prompt matcher changes them by $-0.00122/+0.00134$ relative to the frozen-geometry decoder baseline. The mixed directions keep the claim scoped to complementary evidence rather than uniform improvement.

\noindent\textbf{Sparse hierarchical context.}
We connect aligned fine, middle, and coarse regions through a zero-gated sparse hierarchy, so broader context can be added without perturbing the pretrained fine representation at initialization. A paired 4,000-episode comparison with 10,000 bootstrap resamples measures a macro Dice change of $+0.00018$ and a micro Dice change of $-0.00098$; the latter does not establish the prespecified non-inferiority margin. We therefore report hierarchical interaction as a tested architectural contribution and retain the fine-scale model as the performance-selected primary route, rather than claiming an unsupported universal gain.

\noindent\textbf{Episodic target-versus-rest learning.}
We train the complete prediction pathway episodically across changing target tissues and evaluate it once after patient-separated model selection. On the frozen test set, the final model improves macro/micro Dice over its preceding prompt-conditioned checkpoint by $+0.00126/+0.00285$, improves 7 of 12 tissue targets, and preserves a single inference threshold of 0.5. This experiment directly tests the intended setting: one model, heterogeneous visual prompts, and no class-specific fine-tuning at inference.

\section{Related Work}

\subsection{Interactive and Promptable Segmentation}

Interactive segmentation first encoded user intent through task-specific spatial cues. DEXTR transforms four extreme clicks into an additional image channel for class-agnostic object segmentation \cite{maninis2018dextr}, while FocalClick performs coarse target localization and focused correction to update a mask efficiently from positive and negative clicks \cite{chen2022focalclick}. Foundation models broadened this interface: SAM accepts points, boxes, and masks for promptable segmentation in natural images \cite{kirillov2023segment}, and MedSAM adapts box prompting to medical images at scale \cite{ma2024medsam}. In pathology, WSI-SAM incorporates multi-resolution features into a SAM-based model for whole-slide segmentation \cite{liu2024wsisam}. These methods establish effective prompt-to-mask prediction, but their interaction is primarily resolved within image- or feature-level decoders. CaRePrompt instead uses the prompt to query learned tissue regions whose semantics are jointly determined by cells and cross-scale context.

\subsection{Multiscale Whole-Slide Representation}

WSI analysis commonly separates local feature extraction from broader contextual aggregation. HistoSegNet derives tissue-type maps from patch supervision, class activation maps, and spatial post-processing \cite{chan2019histosegnet}. HIPT organizes nested WSI tiles into a hierarchical self-supervised Transformer to connect local and slide-level representations \cite{chen2022hipt}. More recently, Prov-GigaPath models tens of thousands of tile tokens with a long-context slide encoder and large-scale whole-slide pretraining \cite{xu2024gigapath}. These approaches demonstrate that physical scale and slide context are essential, but they primarily learn representations for predefined tissue labels or slide-level endpoints. Our setting is different: the output class is specified at interaction time, so multiscale information must be routed into a dense target-versus-rest mask conditioned on the current positive and negative evidence.

\subsection{Learned Regions and Cellular Context}

Region-based models reduce dense images to structured prediction units. Superpixel Sampling Networks make superpixel assignment differentiable and allow downstream losses to learn task-specific regions \cite{jampani2018ssn}. Computational pathology adds biologically meaningful entities to this abstraction. HoVer-Net jointly segments and classifies nuclei across multiple tissues \cite{graham2019hover}, and CellViT learns transformer-based cellular representations that transfer across histopathology datasets \cite{horst2024cellvit}. HACT-Net goes further by connecting cell and tissue graphs in a hierarchical representation for breast cancer classification \cite{pati2022hact}. These works motivate both adaptive regions and cell-to-tissue structure. CaRePrompt combines them for a distinct prediction interface: soft tissue regions receive cellular evidence, exchange sparse context across aligned physical scales, and remain directly addressable by signed prompts before their probabilities are projected back to pixels.

\section{Method}

\subsection{Problem Formulation and Overview}

Let $I$ denote an H\&E WSI and $\Omega$ its pixel domain at the reference resolution. An interaction provides a set of positive prompts $\mathcal{P}^{+}$ and negative prompts $\mathcal{P}^{-}$, where each prompt may be a point, box, scribble, or partial region. The prompts specify a visual concept rather than a predefined semantic label. CaRePrompt therefore learns a target-versus-rest mapping
\begin{equation}
    f_{\theta}(I,\mathcal{P}^{+},\mathcal{P}^{-}) = \hat{M},
    \qquad \hat{M}\in[0,1]^{|\Omega|},
\end{equation}
where $\hat{M}$ is the probability mask of the tissue indicated by the current interaction. A single model is trained across changing target classes and can answer a new query without class-specific fine-tuning.

CaRePrompt represents the slide as a hierarchy of learned tissue regions. Spatially aligned views first provide complementary morphology at fine, middle, and coarse scales. A differentiable region encoder converts dense features into soft regions; a cell branch then injects cellular evidence into their tokens. Sparse parent--child interactions connect local morphology to broader tissue organization. Finally, a signed prompt encoder defines the target concept, and a context-aware decoder predicts region probabilities that are projected back to the pixel grid.

\subsection{Multiscale Deep Regionization}

For each sampled location, we extract aligned views $\{I_s\}_{s\in\mathcal{S}}$ with $\mathcal{S}=\{f,m,c\}$ denoting fine, middle, and coarse scales. These views share a center but cover progressively larger physical fields of view. Thus, the fine view preserves cellular morphology and narrow boundaries, while the coarser views capture glands, tissue composition, and global organization without assigning any tissue class to a fixed scale.

At scale $s$, an image encoder produces a dense feature map $F_s\in\mathbb{R}^{H_s\times W_s\times d}$. A region-assignment head predicts $K_s$ soft regions,
\begin{equation}
    A_s(x,k)=
    \frac{\exp(a_s(x,k))}
    {\sum_{j=1}^{K_s}\exp(a_s(x,j))},
\end{equation}
where $A_s(x,k)$ is the membership of pixel $x$ in region $k$. The corresponding region token is obtained by normalized soft pooling,
\begin{equation}
    r_{s,k}=
    \frac{\sum_x A_s(x,k)F_s(x)}
    {\sum_x A_s(x,k)+\varepsilon}.
\end{equation}
Unlike fixed superpixels, both the region boundaries and their representations receive gradients from the downstream task. Region pretraining begins from conventional over-segmentation as weak geometric supervision, while boundary, compactness, balance, and semantic-purity objectives prevent fragmented or collapsed assignments. Subsequent episodic training can then move region boundaries toward structures that are useful for prompt-conditioned segmentation.

\subsection{Cell-Aware Region Representation}

Macroscopic appearance alone can be ambiguous in H\&E images; regions with similar color and texture may contain different cellular populations. We therefore associate each detected cell $i$ with its coordinate $u_i$ and embedding $z_i$. The soft assignment evaluated at $u_i$ determines which regions may receive its evidence. For region $(s,k)$, attention pooling is defined as
\begin{align}
    e_{s,k,i} & =
    \frac{(W_q r_{s,k})^{\top}(W_k z_i)}{\sqrt{d}}, \\
    \alpha_{s,k,i} & =
    \frac{A_s(u_i,k)\exp(e_{s,k,i})}
    {\sum_j A_s(u_j,k)\exp(e_{s,k,j})+\varepsilon}, \\
    c_{s,k} & = \sum_i \alpha_{s,k,i}W_v z_i.
\end{align}
This construction uses the region appearance to select informative cells while retaining cell density and summary statistics as explicit evidence. A learned fusion layer combines $r_{s,k}$, $c_{s,k}$, and these statistics into a cell-aware token $\tilde{r}_{s,k}$. It also represents low-cellularity tissue explicitly rather than treating the absence of cells as missing information.

\subsection{Sparse Hierarchical Region Interaction}

Because the multiscale views are spatially registered, overlap induces parent--child relations between adjacent levels. We retain sparse fine-to-middle and middle-to-coarse edges rather than applying dense attention over all WSI regions. For a parent region $v$, child evidence is aggregated by
\begin{equation}
    m_v = \operatorname{Attn}\!\left(
    \tilde{r}_v,
    \{\tilde{r}_u\}_{u\in\mathcal{C}(v)}
    \right),
\end{equation}
where $\mathcal{C}(v)$ contains its spatially overlapping children. Information is first summarized from fine to middle and then from middle to coarse. The resulting parent and ancestor context is returned to each fine region through a gated residual adapter,
\begin{equation}
    h_{f,k}=\tilde{r}_{f,k}
    +\gamma\,\operatorname{HAttn}\!\left(
    \tilde{r}_{f,k},\tilde{r}_{\pi(k)},\tilde{r}_{\pi^2(k)}
    \right),
\end{equation}
where $\pi(k)$ and $\pi^2(k)$ denote the parent and ancestor regions. The gate $\gamma$ is initialized at zero, allowing hierarchical context to be introduced without perturbing the pretrained fine-scale representation. This sparse exchange lets cellular and boundary evidence constrain large tissue regions while returning broader organization to local predictions.

\subsection{Signed Visual Prompt Encoding}

Prompts are mapped to the current learned regions by their normalized geometry, so their association remains valid as soft assignments evolve. Let $R^{+}$ and $R^{-}$ be the sets of region tokens covered by positive and negative prompts. Separate permutation-invariant set encoders summarize the two prompt polarities,
\begin{equation}
    q^{+}=\operatorname{SetEnc}^{+}(R^{+}),\qquad
    q^{-}=\operatorname{SetEnc}^{-}(R^{-}),
\end{equation}
and form a task token
\begin{equation}
    q=\operatorname{MLP}([q^{+},q^{-},q^{+}-q^{-}]).
\end{equation}
This representation supports variable numbers and types of prompts and preserves variation within each set. For every candidate region, cross-attention to $q^{+}$ and $q^{-}$ produces a prompt-conditioned feature, while signed similarity provides an explicit attraction--repulsion cue. The decoder input is
\begin{equation}
    u_k=[h_{f,k},q,h_k^{\mathrm{prm}},s_k^{+}-s_k^{-},g_k],
\end{equation}
where $h_k^{\mathrm{prm}}$ is the cross-attended prompt feature and $g_k$ contains region geometry. During training, targets and prompt-to-region associations are recomputed from the current assignments rather than inherited from a stale region partition.

\subsection{Context-Aware Mask Decoding}

The mask decoder operates on a sparse graph whose edges connect spatial neighbors and hierarchical relatives. Graph attention propagates local continuity and prompt evidence before a binary head predicts the target probability $p_k=\sigma(o_k)$ for each fine region. The soft region assignments then recover a dense mask:
\begin{equation}
    \hat{M}(x)=\sum_{k=1}^{K_f}A_f(x,k)p_k.
\end{equation}
This direct probability prediction avoids rank cutoffs, classwise area priors, and manually chosen score-to-mask rules. We use a single threshold of $0.5$ at inference. If positive and negative prompts collapse onto the same hard region, the model abstains and requests a refined interaction instead of silently returning an ambiguous mask.

\subsection{Staged Episodic Optimization}

Training follows the dependency structure of the model. We first learn stable multiscale regions, then optimize cell-aware region representations and hierarchical interactions before introducing prompt-conditioned episodes. Each episode samples a WSI crop, a target tissue, one or more positive prompts, optional hard negative prompts, and the corresponding binary mask. Varying the target across episodes teaches the model to infer the task from visual evidence rather than a fixed output channel.

The joint objective is
\begin{equation}
\begin{split}
    \mathcal{L}={}&
    \mathcal{L}_{\mathrm{BCE}}+lambda_d\mathcal{L}_{\mathrm{Dice}}
    +\lambda_b\mathcal{L}_{\mathrm{bd}}
    +\lambda_r\mathcal{L}_{\mathrm{reg}} \\
    &+\lambda_a\mathcal{L}_{\mathrm{assign}}
    +\lambda_t\mathcal{L}_{\mathrm{task}},
\end{split}
\end{equation}
where $\mathcal{L}_{\mathrm{reg}}$ supervises online region targets, $\mathcal{L}_{\mathrm{assign}}$ collects balance, entropy, and compactness regularizers, and $\mathcal{L}_{\mathrm{task}}$ preserves the prompt representation during joint adaptation. The pixel BCE and Dice terms optimize the final interaction objective, while the boundary term encourages precise tissue contours. Joint fine-tuning uses conservative updates to the pretrained image encoder and retains the region regularizers to avoid assignment collapse.

% References and End of Paper
\bibliography{references}

% Check whether the conference requires a reproducibility checklist to be included in the paper.
% If so, you can uncomment the following line and adjust the path to include it.
% \input{ReproducibilityChecklist.tex}

\end{document}
